\chapter*{Introduction}
\addcontentsline{toc}{chapter}{Introduction}

% Contexte du projet, problématique abordée et présentation rapide du sujet
% (RAG sur le Code de la propriété intellectuelle via Legifrance).
% À compléter :
% - Mise en contexte (qu'est-ce qu'un RAG, intérêt sur des textes juridiques)
% - Problématique : pourquoi un RAG sur le Code de la propriété intellectuelle ?
% - Objectifs du projet
% - Plan du rapport

\paragraph{Contexte}
% Décrire brièvement le contexte du cours LO17/AI31 et la notion de RAG.

\paragraph{Problématique}
% Pourquoi appliquer un RAG au Code de la propriété intellectuelle ?
% Difficultés d'accès, volume des textes, besoin de réponses sourcées, etc.

\paragraph{Objectifs}
% Lister les objectifs : ingestion Legifrance, indexation Chroma, génération
% via un LLM, évaluation, interface Streamlit.

\paragraph{Plan du rapport}
% Annonce des parties.
